\documentclass{exam}

\usepackage{amsmath,amssymb,amsfonts}
\usepackage{graphicx}
\usepackage{amsthm}
\usepackage[french]{babel}
\usepackage{enumitem}
\pagestyle{empty}
\graphicspath{ {../Images/} }

\usepackage{geometry}
\geometry{left=1.5cm, right=1.5cm, top=1.5cm, bottom=1.5cm}


%\setlength{\topmargin}{-.5in} \setlength{\textheight}{9.25in}
%\setlength{\oddsidemargin}{0in} \setlength{\textwidth}{6.8in}

\makeatletter
\def\gap#1{%
  \setbox\@tempboxa\hbox{{\color{white}#1}}%
  \@tempdima\wd\@tempboxa
  \rlap{\unhbox\@tempboxa}%
  \hb@xt@\@tempdima{\dotfill}%
}
\makeatother   

\theoremstyle{definition}
\newtheorem{exercice}{Exercice}

\begin{document}
\large

\noindent{\bf Contrôle 3 -- 25min. \hfill Première générale \hfill 20/11/2024}

\vspace{2mm}

\noindent{\makebox[\textwidth]{Nom et prénom:\enspace\hrulefill}}

\vspace{2mm}

\begin{exercice}
Écrire sous la forme $a^n$ avec $a, n \in \mathbb{N}$.
    \begin{enumerate}
        \item $3^3 \times 8^3$
        \item $(6^4)^4$
    \end{enumerate}
    \begin{center}
        \framebox[1\textwidth]{\parbox{5.5in}{\centering \underline{Réponse:}
        \vspace{5em}}}
    \end{center}
\end{exercice}

\begin{exercice}
    Soit $(w_n)$ une suite définie pour tout $n \in \mathbb{N}$ par $w_n = \dfrac{8}{3^n}$.
    
    \begin{enumerate}
        \item Quelle est la nature de cette suite ? Justifier.
        \item Donner sa raison et son premier terme $w_0$.
    \end{enumerate}
    \begin{center}
        \framebox[1\textwidth]{\parbox{5.5in}{\centering \underline{Réponse:}
        \vspace{12em}}}
    \end{center}
\end{exercice}

\begin{exercice}
    On considère la suite définie par $u_n = -n^2 + n + 1$ pour tout $n \in \mathbb{N}$.
    
    \begin{enumerate}
        \item Calculer $u_{n+1} - u_n$ pour tout $\mathbb{N}$.
        \item En déduire les variations de la suite $(u_n)$ sur $\mathbb{N}$.
    \end{enumerate}
    \begin{center}
        \framebox[1\textwidth]{\parbox{5.5in}{\centering \underline{Réponse:}
        \vspace{16em}}}
    \end{center}
\end{exercice}

\begin{exercice}
    Résoudre l'équation suivante en calculant d'abord le discriminant $\Delta$ :
    \[
    x^2 - 2x + 1 = 0
    \]
    \begin{center}
        \framebox[1\textwidth]{\parbox{5.5in}{\centering \underline{Réponse:}
        \vspace{23em}}}
    \end{center}
\end{exercice}

\begin{exercice}
    Factoriser le polynôme suivant en calculant d'abord le discriminant $\Delta$ :
    \[
    -x^2 - 3x + 4
    \]
    \begin{center}
        \framebox[1\textwidth]{\parbox{5.5in}{\centering \underline{Réponse:}
        \vspace{23em}}}
    \end{center}
\end{exercice}

\end{document}